\subsection{The optimal constant in dimension one}
\label{subsec:psh-planar-ideal-distance}

Let $\Omega\subset\mathbb C$ be a pseudoconvex domain, with Euclidean area $\mathrm{d}V$. We use the weighted norms and global ideal sections of \cref{subsec:psh-quantitative-ideals}. In one dimension every multiplier ideal is determined by vanishing orders at points. This makes it possible to combine the selecting weights before applying the one-weight estimate.

\Needspace{9\baselineskip}
\begin{lemma}\label{lem:planar-multiplier-ideal}\index{Multiplier ideal!on a plane domain}
For a genuine subharmonic function $u$ on $\Omega$ and $x\in\Omega$,
\begin{equation}\label{eq:planar-multiplier-ideal}
 \mathcal I(u)_x=(z-x)^{\lfloor\nu(u,x)\rfloor}\mathcal O_{\Omega,x}.
\end{equation}
These ideals form a locally principal ideal sheaf.
\end{lemma}
\begin{proof}
In the normalization of the book, $\mu=\ddc u$ is a positive Radon measure and $a=\mu(\{x\})=\nu(u,x)$. Remove this atom. The distribution $u-a\log|z-x|^2$ has positive complex Hessian and therefore has a genuine subharmonic representative $v$, with $\ddc v(\{x\})=0$. Equality of subharmonic representatives gives
\[
 u=a\log|z-x|^2+v.
\]

For each $t>0$, the function $\mathrm e^{-tv}$ is integrable on some neighborhood of $x$. Indeed, choose a disc $D=D(x,R)\Subset\Omega$ with $M=(\ddc v)(D)<1/t$. Such discs exist because $\ddc v$ has no atom at $x$. On $D$ write
\[
 v(z)=\int_D\log|z-y|^2\,\mathrm{d}(\ddc v)(y)+h(z),
\]
where $h$ is harmonic. The integral is locally integrable by Fubini, has complex Hessian $\ddc v$ on $D$, and hence differs from $v$ by a harmonic distribution. If $M=0$, the assertion follows from harmonicity. For $M>0$, Jensen's inequality for the probability measure $(\ddc v)/M$ gives
\[
 \exp\left(-t\int_D\log|z-y|^2\,\mathrm{d}(\ddc v)(y)\right)
 \leq\frac1M\int_D |z-y|^{-2tM}\,\mathrm{d}(\ddc v)(y).
\]
The integral in $z$ over $D(x,R/2)$ is uniformly finite for $y\in D$, since $tM<1$. The harmonic function is bounded below on this smaller disc, so Tonelli proves the assertion. Truncating the logarithmic kernel justifies Jensen at its poles.

A nonzero holomorphic germ of order $k$ is a unit times $(z-x)^k$. Since $v$ is locally bounded above, weighted integrability requires $k+1>a$, including exclusion of equality by the radial logarithmic divergence. Conversely, if $k+1>a$, choose $b>1$ such that $b(a-k)<1$. The factor $|z-x|^{2k-2a}$ belongs locally to $L^b$, and $\mathrm e^{-v}$ belongs locally to $L^{b/(b-1)}$ by the preceding argument. H\"older proves integrability. Thus the allowed integers are exactly $k\geq\lfloor a\rfloor$.

The points with $\nu(u,y)\geq1$ are locally finite: each contributes an atom of mass at least one to the locally finite measure $\ddc u$. Near any point, therefore, the prescribed ideal is generated by a single power of a coordinate, proving the sheaf assertion.
\end{proof}

The sum in the distance condition consists of global sections. We next justify passing from the corresponding local condition to this global sum.

\Needspace{9\baselineskip}
\begin{lemma}\label{lem:planar-global-ideal-sum}
If $u_1,\ldots,u_r$ are genuine subharmonic functions on $\Omega$, then
\begin{equation}\label{eq:planar-global-ideal-sum}
 \Gamma\left(\Omega,\sum_{i=1}^r\mathcal I(u_i)\right)
       =\sum_{i=1}^r\Gamma(\Omega,\mathcal I(u_i)).
\end{equation}
\end{lemma}
\begin{proof}
Put $J_i=\mathcal I(u_i)$ and $J=\sum_iJ_i$. By \cref{lem:planar-multiplier-ideal}, these sheaves are locally free of rank one. At any point their orders are integers, and the sum has the smallest order. A local generator of an ideal attaining that order is also a generator of $J$. Consequently the sum map $\bigoplus_iJ_i\to J$ is locally split and its kernel $E$ is a holomorphic vector bundle of rank $r-1$. For $r=1$ the assertion is immediate.

Fix $f\in\Gamma(\Omega,J)$. A smooth partition of unity applied to local holomorphic lifts gives a smooth section $s$ of $\bigoplus_iJ_i$ whose image is $f$. Its derivative $\beta=\bar\partial s$ takes values in $E$. It suffices to solve $\bar\partial v=\beta$ in $E$, for then $s-v$ is a holomorphic lift.

We solve this equation with a positive bundle metric. Choose a smooth strictly subharmonic exhaustion $\rho\geq1$ of $\Omega$. The K\"ahler metric
\[
 \omega_*=\mathrm i\,\mathrm{d}z\wedge\mathrm{d}\bar z
                   +\mathrm i\partial\bar\partial\rho^2
\]
is complete: it dominates $2\mathrm i\partial\rho\wedge\bar\partial\rho$, so $\rho$ has bounded gradient, and a curve leaving every compact subset has infinite length.

\Needspace{10\baselineskip}
Equip $E$ with a smooth Hermitian metric $h$. Choose a smooth increasing convex function $\chi$ on $[1,+\infty)$ so that the metric $h_\chi=h\mathrm e^{-\chi(\rho)}$ has
\[
 \mathrm i\Theta_{E,h_\chi}\geq\omega_*\operatorname{Id}_E,
 \qquad \int_\Omega |\beta|_{h_\chi,\omega_*}^2\,\mathrm{d}V_{\omega_*}<+\infty.
\]
These two requirements can be imposed on successive compact sublevels of $\rho$: increase $\chi'$ to dominate the negative curvature relative to the positive Hessian of $\rho$, and increase $\chi$ so that the integral over each shell $\{m\leq\rho<m+1\}$ is at most $2^{-m}$. Increasing the slopes preserves the earlier inequalities.

Use the global canonical frame $\mathrm{d}z$ to identify $E$-valued $(0,q)$-forms with $(E\otimes K_\Omega^{-1})$-valued $(1,q)$-forms. Give $K_\Omega^{-1}$ its flat coordinate metric. Its curvature is zero, and $\omega_*$ dominates the Euclidean metric, so the identified $\beta$ remains $L^2$. The Bochner--Kodaira inequality for compactly supported forms $w$ of type $(1,1)$ gives
\[
 \|\bar\partial^*w\|^2\geq
       \langle[\mathrm i\Theta_{E,h_\chi},\Lambda_{\omega_*}]w,w\rangle
       \geq\|w\|^2.
\]
There are no $(1,2)$-forms. Completeness extends this inequality to the adjoint domain by cutoffs and smoothing. Since $\beta$ is $L^2$, Cauchy inequality bounds $w\mapsto\langle w,\beta\rangle$ by $\|\beta\|\|\bar\partial^*w\|$. Riesz representation gives a distributional solution $v$ with $\bar\partial v=\beta$. Thus $s-v$ is distributionally holomorphic, hence holomorphic in every bundle chart, and supplies the required global lift. This construction imposes no global weighted norm condition on its components.
\end{proof}

\Needspace{17\baselineskip}
\begin{theorem}\label{thm:planar-optimal-ideal-distance}\index{Multiplier ideal!optimal planar distance}
Let $\psi$ and $\phi_1,\ldots,\phi_r<0$ be genuine subharmonic functions on $\Omega$, with $r\geq1$, and choose $p_i>1$. For holomorphic $f$ with $I=\|f\|_{\psi+\sum_i\phi_i}^2<+\infty$, define
\[
 c=\inf\left\{\|g\|_\psi^2:g\in\mathcal O(\Omega),\quad
 f-g\in\sum_{i=1}^r\Gamma(\Omega,\mathcal I(p_i\phi_i+\psi))\right\},
 \qquad s=\sum_{i=1}^r\frac1{p_i}.
\]
Then
\begin{equation}\label{eq:planar-optimal-ideal-distance}
 c\leq\max\{1-s,0\}\,I.
\end{equation}
For every fixed choice of the exponents, this constant is optimal over the stated class of data. If $s\geq1$, the function $g=0$ is admissible.
\end{theorem}
\begin{proof}
Put $J_i=\mathcal I(p_i\phi_i+\psi)$ and $J=\sum_iJ_i$. At a point $x$, write $A=\nu(\psi,x)$ and $B_i=\nu(\phi_i,x)$. By \cref{lem:planar-multiplier-ideal}, the order of $J_x$ is $\min_i\lfloor A+p_iB_i\rfloor$.

Suppose first that $s<1$, and put $P=1/s>1$. The numbers $P/p_i$ sum to one, so
\[
 \min_i p_iB_i\leq\sum_i\frac{P}{p_i}p_iB_i=P\sum_iB_i.
\]
Taking integer parts gives $\mathcal I(P\sum_i\phi_i+\psi)\subseteq J$. By \cref{lem:planar-global-ideal-sum}, every global remainder in this smaller ideal is admissible for $c$. Apply \cref{prop:one-weight-ideal-distance} with $\phi=\sum_i\phi_i$ and $p=P$. It gives $c\leq(1-1/P)I=(1-s)I$.

If $s\geq1$, finite $I$ implies $f\in\Gamma(\Omega,\mathcal I(\psi+\sum_i\phi_i))$. Unless $f$ vanishes identically, let $k$ be its order at $x$. Then $k+1>A+\sum_iB_i$. If its germ did not belong to $J_x$, we would have $k+1\leq A+p_iB_i$ for every $i$. With $L=k+1-A>0$, this would give $\sum_iB_i\geq sL\geq L$, a contradiction. Hence $f\in\Gamma(\Omega,J)$, and \cref{lem:planar-global-ideal-sum} makes $g=0$ admissible. The identically zero function has the same conclusion.

For sharpness when $s<1$, take the unit disc, $\psi=0$, $f=1$, and $\phi_i=p_i^{-1}\log|z|^2$. All $J_i=(z)$, so an admissible $g$ has $g(0)=1$. Orthogonality of disc monomials gives $c=\pi$, attained by $g=1$, whereas $I=\pi/(1-s)$. If $s\geq1$, retain these weights and take $f=z^q$ with an integer $q\geq0$ satisfying $q+1>s$. Then $I=\pi/(q+1-s)>0$ and $c=0$, excluding a negative universal constant.
\end{proof}

For $r\geq2$ on a plane domain, the product constant in \cref{thm:product-ideal-distance} is therefore not optimal. The point-order argument uses the fact that all local ideals are principal; it does not determine the optimal constant for $n\geq2$ and $r>n$.
